% FAC-EN-U0010
% French source: fac_body.tex lines 313--388.
% Authority: PDF physical pp. 11--13; printed pp. 206--208 (no. 11 and no. 12 complete).

\medskip
\phantomsection\label{fac:I:1:no:11}
\textbf{11. Sheaf of germs of homomorphisms from one sheaf to another.}

Let $\mathcal{A}$ be a sheaf of rings, and let $\mathcal{F}$ and $\mathcal{G}$ be two sheaves of $\mathcal{A}$-modules. If $U$ is an open subset of $X$, let $\mathcal{H}_U$ be the group of homomorphisms from $\mathcal{F}(U)$ to $\mathcal{G}(U)$ (we shall also say ``a homomorphism from $\mathcal{F}$ to $\mathcal{G}$ over $U$'' in place of ``a homomorphism from $\mathcal{F}(U)$ to $\mathcal{G}(U)$''). Restriction of a homomorphism defines $\varphi_U^V:\mathcal{H}_V\to\mathcal{H}_U$. The sheaf defined by the system $(\mathcal{H}_U,\varphi_U^V)$ is called the \textit{sheaf of germs of homomorphisms} from $\mathcal{F}$ to $\mathcal{G}$ and is denoted by $\operatorname{Hom}_{\mathcal{A}}(\mathcal{F},\mathcal{G})$. If the $\mathcal{A}_x$ are commutative, $\operatorname{Hom}_{\mathcal{A}}(\mathcal{F},\mathcal{G})$ is a sheaf of $\mathcal{A}$-modules.

An element of $\operatorname{Hom}_{\mathcal{A}}(\mathcal{F},\mathcal{G})_x$, being a germ of a homomorphism from $\mathcal{F}$ to $\mathcal{G}$ in a neighborhood of $x$, unambiguously defines an $\mathcal{A}_x$-homomorphism from $\mathcal{F}_x$ to $\mathcal{G}_x$; hence there is a canonical homomorphism
\[
\rho:\operatorname{Hom}_{\mathcal{A}}(\mathcal{F},\mathcal{G})_x
\longrightarrow
\operatorname{Hom}_{\mathcal{A}_x}(\mathcal{F}_x,\mathcal{G}_x).
\]
In contrast to the operations studied thus far, however, $\rho$ is not in general a bijection; in \hyperref[fac:I:2:no:14]{no.~14} we shall give a sufficient condition for it to be one.

If $\varphi:\mathcal{F}'\to\mathcal{F}$ and $\psi:\mathcal{G}\to\mathcal{G}'$ are homomorphisms, there is an evident homomorphism
\[
\operatorname{Hom}_{\mathcal{A}}(\varphi,\psi):
\operatorname{Hom}_{\mathcal{A}}(\mathcal{F},\mathcal{G})
\longrightarrow
\operatorname{Hom}_{\mathcal{A}}(\mathcal{F}',\mathcal{G}').
\]

Every exact sequence
\[
0\longrightarrow\mathcal{G}\longrightarrow\mathcal{G}'\longrightarrow\mathcal{G}''
\]
gives rise to an exact sequence
\[
0\longrightarrow\operatorname{Hom}_{\mathcal{A}}(\mathcal{F},\mathcal{G})
\longrightarrow\operatorname{Hom}_{\mathcal{A}}(\mathcal{F},\mathcal{G}')
\longrightarrow\operatorname{Hom}_{\mathcal{A}}(\mathcal{F},\mathcal{G}'').
\]
There are also canonical isomorphisms
\[
\operatorname{Hom}_{\mathcal{A}}(\mathcal{A},\mathcal{G})\simeq\mathcal{G},
\]
\[
\operatorname{Hom}_{\mathcal{A}}(\mathcal{F},\mathcal{G}_1+\mathcal{G}_2)
\simeq
\operatorname{Hom}_{\mathcal{A}}(\mathcal{F},\mathcal{G}_1)
+\operatorname{Hom}_{\mathcal{A}}(\mathcal{F},\mathcal{G}_2),
\]
and
\[
\operatorname{Hom}_{\mathcal{A}}(\mathcal{F}_1+\mathcal{F}_2,\mathcal{G})
\simeq
\operatorname{Hom}_{\mathcal{A}}(\mathcal{F}_1,\mathcal{G})
+\operatorname{Hom}_{\mathcal{A}}(\mathcal{F}_2,\mathcal{G}).
\]

\phantomsection\label{fac:I:2}
\section*{\S2. Coherent sheaves of modules}
\addcontentsline{toc}{section}{Chapter I, \S2. Coherent sheaves of modules}

In this section, $X$ denotes a topological space and $\mathcal{A}$ a sheaf of rings on $X$. We suppose that all the $\mathcal{A}_x$, $x\in X$, are commutative and possess an identity element varying continuously with $x$. All sheaves considered through no.~\hyperref[fac:I:2:no:16]{16} are sheaves of $\mathcal{A}$-modules, and all homomorphisms are $\mathcal{A}$-homomorphisms.

\medskip
\phantomsection\label{fac:I:2:no:12}
\textbf{12. Definitions.}

Let $\mathcal{F}$ be a sheaf of $\mathcal{A}$-modules, and let $s_1,\ldots,s_p$ be sections of $\mathcal{F}$ over an open set $U\subset X$. The rule sending every family $f_1,\ldots,f_p$ of elements of $\mathcal{A}_x$ to the element $\sum_{i=1}^{p}f_i\,s_i(x)$ of $\mathcal{F}_x$ gives a homomorphism $\varphi:\mathcal{A}^p\to\mathcal{F}$ defined over $U$ (more precisely, in the notation of \hyperref[fac:I:1:no:4]{no.~4}, $\varphi$ is a homomorphism from $\mathcal{A}^p(U)$ to $\mathcal{F}(U)$). The kernel $\mathcal{R}(s_1,\ldots,s_p)$ of $\varphi$ is a subsheaf of $\mathcal{A}^p$, called the \textit{sheaf of relations} among the $s_i$; the image of $\varphi$ is the subsheaf of $\mathcal{F}$ generated by the $s_i$. Conversely, every homomorphism $\varphi:\mathcal{A}^p\to\mathcal{F}$ defines sections $s_1,\ldots,s_p$ of $\mathcal{F}$ by the formulas
\[
s_1(x)=\varphi_x(1,0,\ldots,0),\qquad\ldots,\qquad
s_p(x)=\varphi_x(0,\ldots,0,1).
\]

\phantomsection\label{fac:I:2:no:12:def:1}
\textsc{Definition 1.} \textit{A sheaf of $\mathcal{A}$-modules $\mathcal{F}$ is said to be of finite type if it is locally generated by finitely many of its sections.}

In other words, for every point $x\in X$ there must exist an open set $U$ containing $x$ and finitely many sections $s_1,\ldots,s_p$ of $\mathcal{F}$ over $U$ such that every element of $\mathcal{F}_y$, $y\in U$, is a linear combination of the $s_i(y)$ with coefficients in $\mathcal{A}_y$. By the preceding discussion, this is equivalent to saying that the restriction of $\mathcal{F}$ to $U$ is isomorphic to a quotient sheaf of some $\mathcal{A}^p$.

\phantomsection\label{fac:I:2:no:12:prop:1}
\textsc{Proposition 1.} \textit{Let $\mathcal{F}$ be a sheaf of finite type. If $s_1,\ldots,s_p$ are sections of $\mathcal{F}$ defined over a neighborhood of a point $x\in X$ and generating $\mathcal{F}_x$, then they generate $\mathcal{F}_y$ for every $y$ sufficiently close to $x$.}

Since $\mathcal{F}$ is of finite type, there are finitely many sections $t_1,\ldots,t_q$ of $\mathcal{F}$ near $x$ which generate $\mathcal{F}_y$ for $y$ sufficiently close to $x$. Since the $s_j(x)$ generate $\mathcal{F}_x$, there exist sections $f_{ij}$ of $\mathcal{A}$ near $x$ such that $t_i(x)=\sum_{j=1}^{p}f_{ij}(x)s_j(x)$. It follows that, for $y$ sufficiently close to $x$,
\[
t_i(y)=\sum_{j=1}^{p}f_{ij}(y)s_j(y),
\]
and hence the $s_j(y)$ generate $\mathcal{F}_y$, q.e.d.

\phantomsection\label{fac:I:2:no:12:def:2}
\textsc{Definition 2.} \textit{A sheaf of $\mathcal{A}$-modules $\mathcal{F}$ is said to be coherent if:}

\begin{enumerate}[label=(\alph*),leftmargin=2.2em]
\item\label{fac:I:2:no:12:def:2:item:a} \textit{$\mathcal{F}$ is of finite type;} 
\item\label{fac:I:2:no:12:def:2:item:b} \textit{if $s_1,\ldots,s_p$ are sections of $\mathcal{F}$ over an open set $U\subset X$, the sheaf of relations among the $s_i$ is of finite type (on the open set $U$).}
\end{enumerate}

Notice the \textit{local} character of \hyperref[fac:I:2:no:12:def:1]{Definitions 1} and \hyperref[fac:I:2:no:12:def:2]{2}.

\phantomsection\label{fac:I:2:no:12:prop:2}
\textsc{Proposition 2.} \textit{Locally, every coherent sheaf is isomorphic to the cokernel of a homomorphism $\varphi:\mathcal{A}^q\to\mathcal{A}^p$.}

This follows immediately from the definitions and from the remarks preceding \hyperref[fac:I:2:no:12:def:1]{Definition 1}.

\phantomsection\label{fac:I:2:no:12:prop:3}
\textsc{Proposition 3.} \textit{Every finite-type subsheaf of a coherent sheaf is coherent.}

Indeed, if a sheaf $\mathcal{F}$ satisfies condition~\hyperref[fac:I:2:no:12:def:2:item:b]{(b)} of \hyperref[fac:I:2:no:12:def:2]{Definition 2}, then every subsheaf of $\mathcal{F}$ plainly satisfies it as well.
